\documentclass[11pt,a4paper,oneside,openany]{stacks-project-book}

% Non-rendering source-note environments retained from the Stacks sources.
\usepackage{verbatim}
\newenvironment{reference}{\comment}{\endcomment}
\newenvironment{slogan}{\comment}{\endcomment}
\newenvironment{history}{\comment}{\endcomment}

% Diagrams and chapter-support packages.
\usepackage[all]{xy}
\xyoption{2cell}
\UseAllTwocells
\usepackage{multicol}
\usepackage{graphicx}

% Japanese mathematical-book typography, inherited from the frozen P12 witness.
\usepackage{fontspec}
\usepackage{xeCJK}
\defaultfontfeatures{Ligatures=TeX}
\setmainfont{texgyretermes-regular.otf}[
  BoldFont=texgyretermes-bold.otf,
  ItalicFont=texgyretermes-italic.otf
]
\setsansfont{texgyreheros-regular.otf}
\setmonofont{lmmono10-regular.otf}[ItalicFont=lmmono10-italic.otf]
\setCJKmainfont{HaranoAjiMincho-Regular.otf}[
  BoldFont=HaranoAjiMincho-Bold.otf,
  ItalicFont=HaranoAjiMincho-Regular.otf
]
\setCJKsansfont{HaranoAjiGothic-Regular.otf}[
  BoldFont=HaranoAjiGothic-Medium.otf
]
\setCJKmonofont{HaranoAjiGothic-Regular.otf}
\xeCJKsetup{PunctStyle=plain,CheckSingle=true}
\XeTeXlinebreaklocale "ja"
\XeTeXlinebreakskip=0pt plus 1pt
\usepackage[
  a4paper,
  inner=26mm,
  outer=24mm,
  top=24mm,
  bottom=24mm,
  headheight=14pt,
  headsep=8mm,
  footskip=12mm
]{geometry}
\AtBeginDocument{\fontsize{10.95pt}{16.5pt}\selectfont}
\setlength{\emergencystretch}{1em}

% Preserve the two frozen legacy tokens using \'etale in math mode while
% rendering their acute accent.  This is a build-only compatibility shim,
% inherited from the bounded P05 precedent; no chapter source is rewritten.
% The LaTeX format rebinds text accents at document start, so install this
% narrow compatibility override after those format hooks have run.
\AtBeginDocument{%
  \let\stacksJaTextAcute\'
  % The final log proves that exactly these two math-mode uses have argument e.
  \renewcommand{\'}[1]{\ifmmode\text{é}\else\stacksJaTextAcute{#1}\fi}%
}

% Japanese structural names.
\renewcommand{\contentsname}{目次}
\renewcommand{\bibname}{参考文献}
\renewcommand{\proofname}{証明}
\renewcommand{\figurename}{図}
\renewcommand{\tablename}{表}
\renewcommand{\appendixname}{付録}

% Japanese numbered chapter heads while retaining the canonical chapter numbers.
\makeatletter
\def\@makechapterhead#1{\global\topskip 7.5pc\relax
  \begingroup
  \fontsize{17}{22}\bfseries\centering
    \ifnum\c@secnumdepth>\m@ne
      \leavevmode \hskip-\leftskip
      \rlap{\vbox to\z@{\vss
          \ifx\chaptername\appendixname
            \centerline{\normalsize\mdseries 付録\enspace\thechapter}
          \else
            \centerline{\normalsize\mdseries 第\thechapter 章}
          \fi
          \vskip 3pc}}\hskip\leftskip\fi
     #1\par \endgroup
  \skip@34\p@ \advance\skip@-\normalbaselineskip
  \vskip\skip@ }
\makeatother

% Current-volume labels resolve locally; absent chapters use permanent tags.
\usepackage{hyperref}
\hypersetup{
  unicode=true,
  pdflang={ja-JP},
  hidelinks,
  pdfborder={0 0 0},
  pdftitle={Stacks Project 日本語版 - 全116章累積版},
  pdfauthor={The Stacks Project contributors; Japanese translation production},
  pdfsubject={Japanese cumulative translation of all 116 chapters of the pinned Stacks Project snapshot}
}
\urlstyle{same}
\input{tagrefs}

% Theorem environments.
\theoremstyle{plain}
\newtheorem{theorem}[subsection]{定理}
\newtheorem{proposition}[subsection]{命題}
\newtheorem{lemma}[subsection]{補題}

\theoremstyle{definition}
\newtheorem{definition}[subsection]{定義}
\newtheorem{example}[subsection]{例}
\newtheorem{exercise}[subsection]{演習}
\newtheorem{situation}[subsection]{状況}

\theoremstyle{remark}
\newtheorem{remark}[subsection]{注}
\newtheorem{remarks}[subsection]{注}

\renewcommand{\contentsname}{目次}
\renewcommand{\proofname}{証明}
\renewcommand{\refname}{参考文献}
\renewcommand{\figurename}{図}
\renewcommand{\tablename}{表}

\numberwithin{equation}{subsection}

% Macros
\def\lim{\mathop{\mathrm{lim}}\nolimits}
\def\colim{\mathop{\mathrm{colim}}\nolimits}
\def\Spec{\mathop{\mathrm{Spec}}}
\def\Hom{\mathop{\mathrm{Hom}}\nolimits}
\def\Ext{\mathop{\mathrm{Ext}}\nolimits}
\def\SheafHom{\mathop{\mathcal{H}\!\mathit{om}}\nolimits}
\def\SheafExt{\mathop{\mathcal{E}\!\mathit{xt}}\nolimits}
\def\Sch{\mathit{Sch}}
\def\Mor{\mathop{\mathrm{Mor}}\nolimits}
\def\Ob{\mathop{\mathrm{Ob}}\nolimits}
\def\Sh{\mathop{\mathit{Sh}}\nolimits}
\def\NL{\mathop{N\!L}\nolimits}
\def\CH{\mathop{\mathrm{CH}}\nolimits}
\def\proetale{{pro\text{-}\acute{e}tale}}
\def\etale{{\acute{e}tale}}
\def\QCoh{\mathit{QCoh}}
\def\Ker{\mathop{\mathrm{Ker}}}
\def\Im{\mathop{\mathrm{Im}}}
\def\Coker{\mathop{\mathrm{Coker}}}
\def\Coim{\mathop{\mathrm{Coim}}}

% Boxtimes
\DeclareMathSymbol{\boxtimes}{\mathbin}{AMSa}{"02}

% Macros for moduli stacks/spaces
\def\QCohstack{\mathcal{QC}\!\mathit{oh}}
\def\Cohstack{\mathcal{C}\!\mathit{oh}}
\def\Spacesstack{\mathcal{S}\!\mathit{paces}}
\def\Quotfunctor{\mathrm{Quot}}
\def\Hilbfunctor{\mathrm{Hilb}}
\def\Curvesstack{\mathcal{C}\!\mathit{urves}}
\def\Polarizedstack{\mathcal{P}\!\mathit{olarized}}
\def\Complexesstack{\mathcal{C}\!\mathit{omplexes}}
\def\Pic{\mathop{\mathrm{Pic}}\nolimits}
\def\Picardstack{\mathcal{P}\!\mathit{ic}}
\def\Picardfunctor{\mathrm{Pic}}
\def\Deformationcategory{\mathcal{D}\!\mathit{ef}}
